\documentclass[11pt]{article}
\usepackage[margin=0.9in]{geometry}
\usepackage{amsmath,amssymb,amsthm,booktabs,mathrsfs}
\usepackage[T1]{fontenc}
\usepackage{lmodern,microtype}
\usepackage[colorlinks=true,linkcolor=blue,urlcolor=blue]{hyperref}
\newtheorem{theorem}{Theorem}[section]
\newtheorem{lemma}[theorem]{Lemma}
\newtheorem{proposition}[theorem]{Proposition}
\newtheorem{corollary}[theorem]{Corollary}
\theoremstyle{remark}
\newtheorem{remark}[theorem]{Remark}
\newcommand{\norm}[1]{\left\|#1\right\|}
\title{Toward a four-dimensional Yang--Mills construction\\
\large Notebook 17: joint fluctuation tails and the one-step Gaussian domain}
\author{Research working document}
\date{17 October 2026}
\begin{document}
\maketitle
\begin{abstract}
We prove joint bad-block bounds for correlated Gaussian fluctuations
from a counting-coordinate covariance bound. The estimates allow
singular observable covariances and overlapping tests. A relative
complex quadratic perturbation has a corresponding local estimate
without a factor growing with the total volume: untested variables
are integrated out before taking absolute values. In four dimensions,
an exact fixed-ratio rescaling supplies the required covariance bound
for both potential values and first differences. The actual reduced
Wilson Hessian at our nonlinear constrained minimizer satisfies the
relative quadratic hypothesis, including its constraint-curvature
term. A separate extraction of the earlier flat effective-action
bounds supplies the current-lattice Gaussian at arbitrary homogeneous
averaging depth. These are inputs to the quantum integration and
iteration required for the existence argument.
\end{abstract}

\section{Which fluctuation estimate is needed}
\label{gs17-target}

Our objective is a nontrivial four-dimensional Yang--Mills quantum
theory. Notebooks 15 and 16 provide the homogeneous finite-regulator
classical minimizer, its ordinary-constraint transfer, local analytic
responses, and the minimized action. The next step integrates the
fluctuations and sums the regions that leave the analytic chart.

The Gaussian roughness calculation of Notebook 16 distinguishes two
different integrations. Integrating every unresolved mode below a
fixed macroscopic scale does not give a uniform smooth-field domain.
Integrating a fixed-ratio step has different coordinates and different
uniform bounds. Both statements must retain the physical pairing:
an inverse precision as a physical Hilbert operator is not its raw
coordinate covariance matrix.

This notebook proves the joint tail estimate needed for prescribed
bad blocks, its complex version, and the fixed-step quadratic
identifications supported by our construction. A good-cell estimate
alone does not supply a regional nonlinear chart on its complement.
The latter construction and the interacting weights are explicit
interfaces for the next step.

The proofs have been checked internally by separate derivations and
audits. The results are finite-regulator theorems on their stated
domains; a continuum existence theorem is the target of their
subsequent integration.

